% ===== Acknowledgements =====
\newpage
\goldrule
\begin{center}
{\large\bfseries\color{rosedeep} Acknowledgements}
\end{center}
\goldrule

\vspace{0.8em}

\noindent
{\cjk 感谢那一位纯粹、自然、坚韧、智慧的"森林女孩"。}

\vspace{0.4em}

\noindent
\textit{My deepest gratitude to the one I call the forest girl---pure, natural, steadfast, and wise. This paper began on a path, with a flower, and with her example of how to see the world. She taught me, without ever saying so, that the most luxurious things cannot be bought---and that the eye which recognises them is the most valuable thing one person can offer another.}

\vspace{1.2em}

\noindent
{\cjk 也愿世界上所有有趣的人，在平凡的路上，都能看见那朵花。}

\vspace{0.4em}

\noindent
\textit{And may all the interesting people in this world, on their ordinary paths, learn to see the flower---the modest, perishable, inexhaustible luxury that grows in the cracks of every ordinary day, free for anyone with the eyes to see it.}

\vspace{1.8em}

\goldrule

\vspace{1.2em}

\noindent{\color{rosedeep}\textbf{Use of Generative AI}}

\vspace{0.4em}

\noindent
The philosophical framework, theoretical claims, and original arguments of this paper are the author's own, developed through extended dialogue with Claude (Anthropic), a large language model that served as an interactive writing and development partner. Claude assisted with drafting, \LaTeX\ typesetting, literature organisation, and formal articulation of arguments whose substance had been established in prior discussion. The author reviewed, directed, and took intellectual responsibility for all content. All original claims---including the coupling-degree definition of relational luxury, the existence-attestation theory, symmetric generability as a necessary condition for just equivalent agency, and subject-embedding misrecognition as the central justice problem---were determined by the author.

\vspace{1.0em}

\goldrule
